% ===== Block B15: 原文 p.140–151 | §6.3 顶点算子、§6.4 KZ 方程的基本解 | 译手 tr-ch1b =====

\subsection{顶点算子}
% [原文 p.140]

本节讨论 [TK] 与 [K] 中发展出来的顶点算子（vertex operator）。顶点算子是构造
KZ 方程解空间的一种便利工具，而这些解空间与共形块（conformal block）空间相对应。
以下我们经常使用记号
\begin{equation}
\kappa = g^* + \ell , \tag{6.20}
\end{equation}
其中 $g^*$ 是单李代数 $\mathfrak{g}$ 的对偶 Coxeter 数（dual Coxeter number），
数 $\ell$ 是 $\hat{\mathfrak{g}}$ 的表示的水平（level）。

% [Tutor Note: $\kappa=g^*+\ell$ 中的 $g^*$ 是数（对偶 Coxeter 数，常记 $h^\vee$），
%  不是李代数；同一句话里的 $\mathfrak{g}$ 才是李代数，请注意正体/花体的区别。
%  $\kappa$ 是本节所有公式的基本常数：仿射代数的 Sugawara 构造给出中心荷
%  $c=\ell\dim\mathfrak{g}/(\ell+g^*)$，而 (V3) 中的共形维数 $\Delta_\lambda$ 也用它作分母，
%  两个 $\kappa$ 是同一个。原文 p.140]

% [原文 p.141]

\begin{definition}[6.9]
设 $\mathbf{v}=\binom{\lambda}{\mu_2\ \mu_1}$。类型为 $\mathbf{v}$ 的顶点算子
$\Phi(z)$ 是 $X_1=\C^*$ 上的多值
$\Hom_\C(V_\lambda\otimes\cH_{\mu_1},\hat{\cH}_{\mu_2})$ 值全纯函数，
使得对 $u\in V_\lambda$，若置
\[
\Phi(u;z)(v)=\Phi(z)(u\otimes v),\qquad v\in\cH_{\mu_1},
\]
则 $\Phi(u;z)$ 满足以下条件：
\begin{enumerate}[label=(V\arabic*)]
\item $\Phi(u;z)$ 关于 $u\in V_\lambda$ 是线性的。
\item 对任意 $X\in\mathfrak{g}$ 有 $[X(m),\Phi(u;z)]=z^m\Phi(Xu,z)$。
\item $[L_m,\Phi(u;z)]=z^m\left(z\dfrac{\dd}{\dd z}+(m+1)\Delta_\lambda\right)
      \Phi(u;z)$，其中
\[
\Delta_\lambda=\frac{(\lambda,\lambda+2\rho)}{2\kappa}.
\]
\end{enumerate}
\end{definition}

数 $\Delta_\lambda$ 称为顶点算子 $\Phi(z)$ 的共形维数（conformal dimension）。

顶点 $\mathbf{v}=\binom{\lambda}{\mu_1\ \mu_2}$ 可象征性地画成
\figomitn{顶点 $\mathbf{v}$ 的三叉图：$\lambda$ 在上，箭头指向顶点；
$\mu_2$ 在左，箭头背离顶点；$\mu_1$ 在右，箭头指向顶点。}

% [Tutor Note: 图上 $\lambda$ 的箭头朝下（进入顶点）、$\mu_2$ 的箭头朝左（离开顶点）、
%  $\mu_1$ 的箭头朝左（进入顶点）；这与 (V2)、(V3) 中 $X(m)$、$L_m$ 只作用在
%  $V_\lambda\otimes\cH_{\mu_1}$ 一侧相符。注意原文此处把下排印作 $(\mu_1\ \mu_2)$，
%  与定义 6.9 的 $(\mu_2\ \mu_1)$ 次序相反，二者必有一处笔误；本译本照原文排，不改。
%  原文 p.141]

\begin{proposition}[6.10]
(1) 类型为 $\mathbf{v}=\binom{\lambda}{\mu_1\ \mu_2}$ 的任意顶点算子 $\Phi(u;z)$
都有 Laurent 展开
\begin{equation}
\Phi(u;z)=\sum_{m\in\Z}\Phi_m(u)z^{-m-\hat{\Delta}_{\mathbf{v}}} \tag{6.21}
\end{equation}
其中
\begin{equation}
\hat{\Delta}_{\mathbf{v}}=\Delta_\lambda+\Delta_{\mu_1}-\Delta_{\mu_2}, \tag{6.22}
\end{equation}
并且每个 $\Phi_m(u)$ 满足
\[
[L_0,\Phi_m(u)]=(\Delta_{\mu_2}-\Delta_{\mu_1}-\Delta_\lambda)\Phi_m(u),
\qquad m\in\Z .
\]
这意味着
\[
\Phi_m(u):\cH_{\mu_1}(d)\to\cH_{\mu_2}(d-m)
\]
对任意 $d\ge0$ 成立。
(2) 对每个 $u\in V_\lambda$ 有
\[
[X(m),\Phi_{m'}(u)]=[X,\Phi_{m+m'}(u)]=\Phi_{m+m'}(Xu),
\qquad X\in\mathfrak{g},\ m,m'\in\Z,
\]
以及
\[
[L_m,\Phi_{m'}(u)]=\{(m+1)\Delta_\lambda-m-m'-\hat{\Delta}(\mathbf{v})\}
\Phi_{m+m'}(u),\qquad m,m'\in\Z .
\]
\end{proposition}

% [Tutor Note: $\hat{\Delta}_{\mathbf{v}}=\Delta_\lambda+\Delta_{\mu_1}-\Delta_{\mu_2}$
%  的取法正是为了保证 (6.21) 是 $z$ 的整数幂级数：由 (V3) 知 $L_0$ 的平移量为
%  $\Delta_\lambda$，而 $\Phi_m$ 又把 $\cH_{\mu_1}$ 的 $L_0$-特征值 $d$ 搬到
%  $d-m+(\Delta_{\mu_2}-\Delta_{\mu_1})$，两项相消后分量落在同一权空间上。
%  原文 p.141]

\noindent\textit{证明.}\quad 我们把 $\Phi(u,z)$ 分解成和
\[
\Phi(u;z)=\sum_{n\in\Z}\Phi_n(u;z),
\]
其中
\[
\Phi_n(u;z):\cH_{\mu_1}(d)\to\cH_{\mu_2}(d-n)
\]
对任意 $d\ge0$ 成立。于是由 (2.21) 有
\[
[L_0,\Phi_n(u;z)]=(\Delta_{\mu_2}-\Delta_{\mu_1}-n)\Phi_n(u;z).
\]
% [原文 p.142]

另一方面，由上面的条件 (V3) 有
\[
[L_0,\Phi_n(u;z)]=\left(z\dfrac{\dd}{\dd z}+\Delta_\lambda\right)\Phi_n(u;z).
\]
由 (6.22) 有
\[
z\dfrac{\dd}{\dd z}\Phi_n(u;z)=-(\hat{\Delta}_{\mathbf{v}}+n)\Phi_n(u;z).
\]
于是我们得出
\[
\Phi_n(u:z)=\Phi_n(u)z^{-n-\hat{\Delta}_{\mathbf{v}}},
\]
从而有展开
\[
\Phi(u;z)=\sum_{n\in\Z}\Phi_n(u)z^{-n-\hat{\Delta}_{\mathbf{v}}}.
\]
\hfill$\square$

% [Tutor Note: 上式中 $\Phi_n(u:z)$ 的冒号是原文的排版笔误（应为分号），照排不改。
%  另外第一步把 $\Phi(u;z)$ 按「权的平移量」分解成 $\sum_n\Phi_n$，这一步的合法性
%  来自 $\cH_{\mu_1}=\bigoplus_{d\ge0}\cH_{\mu_1}(d)$ 的权空间分解与
%  $\Phi$ 的 (V2)、(V3) 两条性质；原文未展开，此处补足。原文 p.142]

$\Phi(u;z)$ 的主枝（principal branch）这样选取：使 $z^{-\hat{\Delta}_{\mathbf{v}}}$
在 $\R_+=\{z\in\R\mid z>0\}$ 上取正值。

% [Tutor Note: 由于指数 $-\hat{\Delta}_{\mathbf{v}}$ 一般不是整数，$z^{-\hat{\Delta}_{\mathbf{v}}}$
%  是多值的；「主枝」即沿 $\R_+$ 取正实值的那个分支。这解释了定义 6.9 中
%  「多值全纯函数」的说法，也说明下文 (6.23) 只能作为多值截面来读。
%  原文 p.142]

若把 $\hat{\cH}_{\mu_2}$ 按
\[
wX(m)=-X(-m)w,\qquad w\in\hat{\cH}_{\mu_2}
\]
看作右 $\hat{\mathfrak{g}}$-模，则由推论 2.18，它典范同构于右
$\hat{\mathfrak{g}}$-模 $\hat{\cH}^{\dagger}_{\mu_2}$。因此，
$\Hom_\C(V_\lambda\otimes\cH_{\mu_1},\hat{\cH}_{\mu_2})$ 值函数 $\Phi(z)$
可以看作 $V^{\dagger}_\lambda\otimes\cH^{\dagger}_{\mu_1}
\widehat{\otimes}\cH^{\dagger}_{\mu_2}$ 值全纯函数 $\langle\Phi(z)|$。
对 $u\in V_\lambda$、$v\in\cH_{\mu_1}$ 与 $w\in\cH_{\mu_2}$，由引理 2.17
与推论 2.18 有
\begin{equation}
\langle\Phi(z)|u\otimes v\otimes w\rangle=(\Phi(u;z)(v)|w). \tag{6.23}
\end{equation}

% [Tutor Note: $\dagger$ 表示对偶（contragredient）模：$V^\dagger_\lambda$ 上的
%  $\hat{\mathfrak{g}}$-作用取负转置，故 $\langle\Phi(z)|$ 是 $\Phi(z)$ 的转置。
%  右模结构 $wX(m)=-X(-m)w$ 与推论 2.18 的配对一起，把
%  $\Hom_\C(V_\lambda\otimes\cH_{\mu_1},\hat{\cH}_{\mu_2})$ 变成三张量积
%  $V^\dagger_\lambda\otimes\cH^\dagger_{\mu_1}\widehat{\otimes}\cH^\dagger_{\mu_2}$
%  （中间的 $\widehat{\otimes}$ 是完备化张量积，因为 $\hat{\cH}_{\mu_2}$ 是完备化的）。
%  原文 p.142]

定义全纯映射 $s:X_1=\C^*\ni z\mapsto((1:z),z)\in\PP^1\times X_1$，并置
\[
\mathfrak{F}=\{\pi:\PP^1\times X_1\to X_1;0,s,\infty;\zeta,\zeta-z,\eta\}
\]
为带形式邻域的三点 $\PP^1$ 族，其中 $0$（resp. $\infty$）是全纯截面
$X_1\ni z\mapsto((1:0),z)\in\PP^1\times X_1$
（resp. $X_1\ni z\mapsto((0:1),z)\in\PP^1\times X_1$），
而 $\zeta$ 是 $\PP^1\setminus\{\infty\}$ 上的坐标。

% [Tutor Note: 这里 $0=(1:0)$、$\infty=(0:1)$、$s(z)=(1:z)$ 是 $\PP^1$ 的三个点，
%  而 $\zeta$、$\zeta-z$、$\eta$ 分别是这三个点处的形式坐标：在 $0$ 处用 $\zeta$，
%  在 $s(z)$ 处用 $\zeta-z$（它在 $s(z)$ 处为零），在 $\infty$ 处用 $\eta=1/\zeta$。
%  这正是 §3 中「带形式邻域的三点射影直线族」的特例，$\mathfrak{F}$ 是共形块的
%  定义域。原文 p.142]

\begin{theorem}[6.11]
在典范同构
\[
\Hom_\C(V_\lambda\otimes\cH_{\mu_1},\hat{\cH}_{\mu_2})
\simeq V^{\dagger}_\lambda\otimes\cH^{\dagger}_{\mu_1}
\widehat{\otimes}\cH^{\dagger}_{\mu_2}
\]
之下，类型为 $\mathbf{v}=\binom{\lambda}{\mu_1\ \mu_2}$ 的顶点算子 $\Phi(z)$
可以看作 $X_1=\C^*$ 上共形块层
$\cV^{\dagger}_{\lambda\mu_1\mu_2^{\dagger}}(\mathfrak{F})$
的多值全纯截面（由定理 3.18）。
\end{theorem}

% [Tutor Note: 这就是本章的「态–场对应」（state–field correspondence）：
%  顶点算子 $\Phi(z)$ 与共形块层的多值全纯截面是同一对象的两种读法——
%  前者按模式 $\Phi_m(u)$ 展开，后者按三点共形块 $(\lambda;\mu_1,\mu_2)$ 分解；
%  定理的证明正是把 (V1)–(V3) 逐步翻译成共形块的零化条件。原文 p.142]

\noindent\textit{证明.}\quad 只需证明对每个 $z\in X_1$，$\langle\Phi(z)|$ 都是
$\cV^{\dagger}_{\lambda\mu_1\mu_2^{\dagger}}(\mathfrak{F})$ 的元素。由定理 3.18，
共形块 $\cV^{\dagger}_{\lambda\mu_1\mu_2^{\dagger}}(\mathfrak{F})$ 由
$V^{\dagger}_\lambda\otimes\cH^{\dagger}_{\mu_1}
\widehat{\otimes}\cH^{\dagger}_{\mu_2}$ 中被
$\mathfrak{g}\otimes H^0(\PP^1,\cO_{\PP^1}(*(0+\infty)))$ 的作用零化掉的
子空间给出。注意

% [原文 p.143]

注意 $H^0(\PP^1,\cO_{\PP^1}(*(0+\infty)))$ 的一组基由
$\{\zeta^m\}_{m\in\Z}$ 给出。对 $f=\zeta^m$ 与 $X\in\mathfrak{g}$ 有
\[
\begin{aligned}
&\langle\Phi(z)X[f]|u\otimes v\otimes w\rangle\\
&\quad=\langle\Phi(z)|z^mXu\otimes v\otimes w\rangle
 +\langle\Phi(z)|u\otimes X(m)v\otimes w\rangle
 +\langle\Phi(z)|u\otimes v\otimes X(-m)w\rangle\\
&\quad=z^m(\Phi(Xu;z)(v)|w)+(\Phi(u;z)(X(m)v)|w)
 +(\Phi(u;z)(v)|X(-m)w)\\
&\quad=(z^m\Phi(Xu;z)(v)|w)+(\Phi(u;z)(X(m)v)|w)
 -(X(m)\Phi(u;z)(v)|w)\\
&\quad=z^m\Phi(Xu;z)(v\otimes w)-[X(m),\Phi(u;z)](v\otimes w)=0,
\end{aligned}
\]
这里用了条件 (V2)。于是条件 (V2) 不是别的，正是共形块上的规范条件
（gauge condition）。这就证明了定理。\hfill$\square$

% [Tutor Note: 第三行到第四行的关键是 Hermitian 形式与右模结构的相容性：
%  由 $wX(m)=-X(-m)w$ 得 $(\Phi(u;z)(v)|X(-m)w)=-(X(m)\Phi(u;z)(v)|w)$，
%  即把 $X(-m)$ 从右搬到左；第四行再代入 (V2) 便得零。
%  「规范条件」指共形块上 $\mathfrak{g}\otimes H^0$ 的作用为零。
%  原文 p.143]

现在还需说明：若 $\cV^{\dagger}_{\lambda\mu_1\mu_2^{\dagger}}(\mathfrak{F}(z))
\ne\{0\}$，则条件 (V3) 给出非零的顶点算子 $\Phi(z)$。

\begin{proposition}[6.12]
类型为 $\mathbf{v}=\binom{\lambda}{\mu_2\ \mu_1}$ 的顶点算子 $\Phi(u;z)$
通过
\begin{equation}
\varphi(u,v,w)=\langle\Phi(u;z)z^{\hat{\Delta}_{\mathbf{v}}}(v)|w\rangle
\big|_{z=0}=\langle\Phi_0(u)(v),w\rangle \tag{6.24}
\end{equation}
决定一个三线性型
$\varphi\in\Hom_{\mathfrak{g}}(V^{*}_{\mu_2}\otimes V_\lambda\otimes V_{\mu_1},\C)$，
其中 $u\in V_\lambda$，$v\in V_{\mu_1}$，$w\in V_{\mu_2}$。
这个三线性型 $\varphi$ 称为顶点算子 $\Phi(u;z)$ 的初始项（initial term）。

反之，顶点算子 $\Phi(u;z)$ 由它的初始项
$\varphi\in\Hom_{\mathfrak{g}}(V^{*}_{\mu_2}\otimes V_\lambda\otimes V_{\mu_1},\C)
=\Hom_{\mathfrak{g}}(V_\lambda\otimes V_{\mu_1},V_{\mu_2})$
按 (6.24) 唯一决定。
\end{proposition}

% [Tutor Note: (6.24) 中乘 $z^{\hat{\Delta}_{\mathbf{v}}}$ 是为了消去 (6.21)
%  的非整数幂：$\Phi(u;z)z^{\hat{\Delta}_{\mathbf{v}}}=\sum_m\Phi_m(u)z^{-m}$
%  是 $z$ 的 Laurent 级数，令 $z=0$ 只剩 $m=0$ 的 $\Phi_0(u)$；故初始项
%  $\varphi$ 与顶点算子 $\Phi$ 互相唯一决定。原文 p.143]

以下把以 $\varphi$ 为初始项的顶点算子 $\Phi$ 常记作 $\Phi_\varphi$。

命题的证明与命题 6.1 的证明本质上相同，我们只需再考虑附加条件 (V3)。

\begin{proposition}[6.13]
类型为 $\mathbf{v}=\binom{\lambda}{\mu_2\ \mu_1}$ 的顶点算子 $\Phi(z)$
满足下列等式：
\[
\begin{aligned}
&X(m)\Phi(z)=0,\qquad m\ge1,\ X\in\mathfrak{g},\\
&X\Phi(u;z)=\Phi(Xu;z),\qquad X\in\mathfrak{g},\ u\in V_\lambda,\\
&L_m\Phi(z)=0,\qquad m\ge1,\\
&L_0\Phi(z)=\Delta_\lambda\Phi(z),\\
&L_{-1}\Phi(z)=\dfrac{\partial}{\partial z}\Phi(z),\\
&X_\theta(-1)^{\ell-(\lambda,\theta)+1}\Phi(|\lambda\rangle;z)=0.
\end{aligned}
\]
\end{proposition}

% [Tutor Note: 前两式说明 $\Phi(z)$ 是 $\mathfrak{g}$-模同态（$X(m)\Phi=0$ 对
%  $m\ge1$ 即 $\Phi$ 与「消灭算符」对易到零），第三、四式说明它在
%  Virasoro 代数下是最高权向量（$L_m\Phi=0$，$m\ge1$；$L_0\Phi=\Delta_\lambda\Phi$），
%  第五式是平移不变性（$L_{-1}=\partial_z$）。这四式合起来就是「$\Phi$ 是
%  权 $\Delta_\lambda$ 的初级场」。原文 p.143]
% [Tutor Note: 末式是水平 $\ell$ 的可积性条件（null vector 条件）：$X_\theta$ 是
%  与最高根 $\theta$ 对应的根向量，$(\lambda,\theta)$ 是 $\lambda$ 在 $\theta$
%  方向的分量。对水平 $\ell$ 的可积模有 $(\lambda,\theta)\le\ell$，故指数
%  $\ell-(\lambda,\theta)+1\ge1$；该式表示 $\Phi(|\lambda\rangle;z)$ 被
%  $X_\theta(-1)$ 作用足够多次后归零，即它是水平 $\ell$ 的「截断」态。
%  原文 p.143]

\begin{definition}[6.14]
对顶点 $\mathbf{v}=\binom{\lambda}{\mu_2\ \mu_1}$，置
\begin{equation}
\cV(\mathbf{v})=W^{\lambda}_{\bar{\chi}}, \tag{6.25}
\end{equation}
其中 $W^{\lambda}_{\bar{\chi}}$ 由推论 6.2 的 3) 定义，即
$\varphi\in W^{\lambda}_{\bar{\chi}}$ 由下述条件刻画：
\[
\operatorname{proj}_{W_{j_2}}\circ\varphi(W_j\otimes W_{j_1})=0
\]
对 $V_\lambda$、$V_{\mu_1}$、$V_{\mu_2}$ 的任意单 $\mathfrak{t}_\theta$-子模
$W_j$、$W_{j_1}$、$W_{j_2}$ 成立，只要 $j+j_1+j_2>\ell$。
\end{definition}

% [Tutor Note: $\mathfrak{t}_\theta$ 是由最高根 $\theta$ 对应的
%  $\mathfrak{sl}_2$-三元组 $(X_\theta,X_{-\theta},\theta^\vee)$ 张成的子代数；
%  $V_\lambda$ 作为 $\mathfrak{t}_\theta$-模分解成 $W_j$（$j$ 为权）。
%  条件 $j+j_1+j_2>\ell$ 是水平 $\ell$ 的可积性截断：三权之和超过水平时
%  对应分量必须为零。这是推论 6.2 中 $W^{\lambda}_{\bar{\chi}}$ 的定义。
%  原文 p.143]

% [原文 p.144]

\figomitn{顶点 $\mathbf{v}=\binom{\lambda}{\mu_2\ \mu_1}$ 的三叉图：
$\lambda$ 在上，箭头指向顶点；$\mu_2$ 在左，箭头背离顶点；$\mu_1$ 在右，
箭头指向顶点。}

于是由推论 6.2 得到下述定理。

\begin{theorem}[6.15]
类型为 $\mathbf{v}=\binom{\lambda}{\mu_2\ \mu_1}$ 的顶点算子的初始项空间，
与向量空间 $\cV(\mathbf{v})$ 等同。
\end{theorem}

置
\[
N^{\mu_2}_{\lambda\mu_1}=\dim_\C\cV^{\dagger}_{\lambda\mu_1\mu_2^{\dagger}}
(\mathfrak{X}),
\]
其中
\[
\mathfrak{X}=(\PP^1;0,1,\infty;\zeta,\zeta-1,\eta),\qquad \eta=1/\zeta .
\]
则由推论 6.2 的 3) 有
\begin{equation}
\dim_\C\cV^{\dagger}_{\lambda\mu_1\mu_2^{\dagger}}(\mathfrak{X})
=\dim_\C\cV(\mathbf{v}), \tag{6.26}
\end{equation}
其中 $\mathbf{v}=\binom{\lambda}{\mu_2\ \mu_1}$。

设 $A_j:\cH_{\mu_{j-1}}\to\hat{\cH}_{\mu_j}$，$j=1,\dots,n$，是一列算子。
注意带右 $\hat{\mathfrak{g}}$-作用的 $\hat{\cH}_{\mu_j}$ 可与带左
$\hat{\mathfrak{g}}$-作用的 $\cH^{\dagger}_{\mu_j}$ 等同。我们固定
$\cH_{\mu_j}$ 与 $\cH^{\dagger}_{\mu_j}$ 的对偶基
\[
\{u^{(j)}_{d,1},\dots,u^{(j)}_{d,m_d}\}\subset\cH_{\mu_j},\qquad
\{u^{(j)*}_{d,1},\dots,u^{(j)*}_{d,m_d}\}\subset\cH^{\dagger}_{\mu_j},
\]
并对任意 $\langle v|\in\cH^{\dagger}_{\mu_n}$ 与 $|u\rangle\in\cH_{\mu_0}$，
形式地定义 $\langle v|A_n\cdots A_1|u\rangle$ 为
\[
\begin{aligned}
\langle v|A_n\cdots A_1|u\rangle
={}&\sum_{d_1,\dots,d_{n-1}\ge0}\sum_{j_1=1}^{m_{d_1}}\cdots
 \sum_{j_{n-1}=1}^{m_{d_{n-1}}}
 \langle v|A_n|u^{(n-1)}_{d_{n-1},j_{n-1}}\rangle\\
&\times\langle u^{(n-1)*}_{d_{n-1},j_{n-1}}|A_{n-1}
 |u^{(n-2)}_{d_{n-2},j_{n-2}}\rangle\cdots
 \langle u^{(1)*}_{d_1,j_1}|A_1|u\rangle .
\end{aligned}
\]

\begin{definition}[6.16]
称算子序列 $\{A_n,A_{n-1},\dots,A_1\}$ 是可复合的（composable），
若对任意 $v_1\in\cH^{\dagger}_{\mu_n}$ 与 $u_1\in\cH_{\mu_1}$，级数
\[
\begin{aligned}
\sum_{d_1,\dots,d_{m-1}\ge0}\Bigg|&
\sum_{j_1=1}^{m_{d_1}}\cdots\sum_{j_{n-1}=1}^{m_{d_{n-1}}}
\langle v|A_n|u^{(n-1)}_{d_{n-1},j_{n-1}}\rangle\\
&\times\langle u^{(n-2)*}_{d_{n-2},j_{n-2}}|A_{n-3}
 |u^{(n-3)}_{d_{n-3},j_{n-3}}\rangle\cdots
 \langle u^{(1)*}_{d_1,j_1}|A_1|u_1\rangle\Bigg|
\end{aligned}
\]
总收敛。
\end{definition}

若 $\{A_n,A_{n-1},\dots,A_1\}$ 可复合，则
$\langle v_1|A_{n_m}\cdots A_{n_1}|u\rangle$ 有确定的值。

把上述定义用于顶点算子序列
$\{\Phi_{\mathbf{v}_n},\Phi_{\mathbf{v}_{n-1}},\dots,\Phi_{\mathbf{v}_1}\}$，其中
\[
\mathbf{v}_n=\binom{\lambda_n}{0\ \mu_{n-1}},\quad
\mathbf{v}_{n-1}=\binom{\lambda_{n-1}}{\mu_{n-1}\ \mu_{n-2}},\
\dots,\
\mathbf{v}_2=\binom{\lambda_2}{\mu_2\ \mu_1},\quad
\mathbf{v}_1=\binom{\lambda_1}{\mu_1\ 0}.
\]
我们把诸顶点用下述符号图形表示。

\figomitn{顶点算子链的符号图：一条水平线（表示 $\PP^1$）上立着
$\lambda_n,\lambda_{n-1},\dots,\lambda_1$ 诸顶点，相邻顶点之间标 $\mu_{n-1},
\mu_{n-2},\dots,\mu_1$，线段两端标 $0$。}

% [原文 p.145]

取 $w_j\in\cH_{\lambda_j}$，$j=n,n-2,\dots,1$，置
$A_j(w_j)=\Phi_{\mathbf{v}_j}(w_j;z_j)$，并定义
\[
\Phi_{\bar{\mathbf{p}}}(z_n,z_{n-1},\dots,z_1)(w_n,w_{n-1},\dots,w_1)
=\langle0|\Phi_{\mathbf{v}_n}(w_n;z_n)\cdots\Phi_{\mathbf{v}_1}(w_1;z_1)|0\rangle,
\]
其中
\[
\bar{\mathbf{p}}=(\mathbf{v}_n,\mathbf{v}_{n-1},\dots,\mathbf{v}_1).
\]
不难证明序列
\[
\{\Phi_{\mathbf{v}_n}(w_n),\Phi_{\mathbf{v}_{n-1}}(w_{n-1}),\dots,
\Phi_{\mathbf{v}_1}(w_1)\}
\]
是可复合的，从而 $\Phi_{\bar{\mathbf{p}}}(z_n,z_{n-1},\dots,z_1)$ 是
$\Hom_\C(V_{\lambda_n}\otimes V_{\lambda_{n-1}}\otimes\cdots\otimes
V_{\lambda_1},\C)$ 中的元素。

$\Phi_{\bar{\mathbf{p}}}(z_n,z_{n-1},\dots,z_1)$ 在 [TK] 与 [K] 中称为
$n$-点函数（$n$-point function），常记作
\[
\langle\Phi_{\mathbf{v}_n}(z_n)\cdots\Phi_{\mathbf{v}_1}(z_1)\rangle .
\]
于是可以证明 $n$-点函数是 KZ 方程的解。更确切地说，有下述定理
（见 [TK]、[K]）。

\begin{theorem}[6.17]
$n$-点函数 $\Phi_{\bar{\mathbf{p}}}(z_n,z_{n-1},\dots,z_1)$ 满足下列方程。
\begin{enumerate}[label=\arabic*.]
\item （规范不变性）对任意 $X\in\mathfrak{g}$ 有
\[
\sum_{j=1}^n\rho_j(X)\Phi_{\bar{\mathbf{p}}}(z_n,z_{n-1},\dots,z_1)=0 .
\]
\item 对每个 $j$ 有
\[
\left(\kappa\frac{\partial}{\partial z_j}
-\sum_{k\ne j}\frac{\Omega_{jk}}{z_j-z_k}\right)
\Phi_{\bar{\mathbf{p}}}(z_n,z_{n-1},\dots,z_1)=0 .
\]
\item 对每个 $j$（$1\le j\le n$）与 $w_k\in V_{\lambda_k}$，置
$L_j=\ell-(\lambda_j,\theta)+1$。则有
\[
\begin{aligned}
&\sum_{|\bar{m}_j|=L_j}\binom{L_j}{\bar{m}_j}
 \prod_{k\ne j}(z_k-z_j)^{-m_k}
 \langle\Phi_{\mathbf{v}_n}(X_\theta^{m_n}u_n;z_n)\cdots
 \Phi_{\mathbf{v}_j}(|\lambda_j\rangle;z_j)\cdots\\
&\qquad\qquad\cdots\Phi_{\mathbf{v}_1}(X_\theta^{m_1}u_1;z_1)\rangle=0,
\end{aligned}
\]
其中 $\bar{m}_j=(m_n,\dots,\widehat{m_j},\dots,m_1)\in(\Z_{\ge0})^{n-1}$，
$|\bar{m}_j|=\sum_{k\ne j}m_k$，而 $\binom{L_j}{\bar{m}_j}$ 是多重组合系数。
\end{enumerate}
\end{theorem}

% [Tutor Note: 定理 6.17 的第 2 条就是 KZ 方程：$\Omega_{jk}$ 是 $\mathfrak{g}$
%  在 $V_{\lambda_j}\otimes V_{\lambda_k}$ 上作用的不变张量（Casimir 的一半），
%  分母 $z_j-z_k$ 是标准的「长程相互作用」形式。第 1 条（规范不变性）保证
%  解落在 $\Hom_{\mathfrak{g}}$ 中，第 3 条是水平 $\ell$ 的可积性（退化）条件，
%  它把解空间限制成有限维的；三条合起来正是 (6.27)、(6.28) 的内容。
%  原文 p.145]

\subsection{KZ 方程的基本解}

{\emergencystretch=5.2pt 设 $\mathbf{A}$ 是标号序列 $\{\lambda_n,\dots,\lambda_1\}$，$\lambda_j\in P_\ell$。
则关于 $\Hom_{\mathfrak{g}}(V_{\lambda_n}\otimes\cdots\otimes V_{\lambda_1},\C)$
值函数 $\Phi(z_n,\dots,z_1)$ 的偏微分方程组\par}
\begin{equation}
\left(\kappa\frac{\partial}{\partial z_i}
-\sum_{k\ne i}\frac{\Omega_{ik}}{z_i-z_k}\right)\Phi(z_n,\dots,z_1)=0,
\qquad i=1,\dots,n \tag{6.27}
\end{equation}
称为 KZ 方程。条件
\begin{equation}
\begin{gathered}
\sum_{\bar{m}_i=L_i}\binom{L_i}{\bar{m}_i}
\prod_{k\ne i}(z_k-z_i)^{-m_k}
\langle\Phi(z_n,\dots,z_1)(X_\theta^{m_n}u_n,\dots,|\lambda_j\rangle,\dots,
X_\theta^{m_1}u_1)\rangle=0,\\
i=1,\dots,n,\qquad \forall u_j\in V_\lambda ,
\end{gathered}
\tag{6.28}
\end{equation}

% [原文 p.146]

其中
\[
\bar{m}_i=(m_n,\dots,\widehat{m_i},\dots,m_1)\in(\Z_{\ge0})^{n-1},\qquad
|\bar{m}_i|=\sum_{k\ne i}m_k ,
\]
称为 $\ell C$-约束（$\ell C$-constraints）。

我们已经说明，$n$-点函数是带 $\ell C$-约束 (6.28) 的 KZ 方程 (6.27) 的解
（定理 6.15）。下面将证明：带 $\ell C$-约束的 KZ 方程的解空间由 $n$-点函数张成。
首先注意算子 $\Omega_{ij}$ 满足下列关系：当 $i,k,m,n$ 互不相同时
$[\Omega_{ik},\Omega_{mn}]=0$；当 $i,k,m$ 互不相同时
$[\Omega_{im},\Omega_{ik}+\Omega_{km}]=0$。

为简化记号，置
\begin{equation}
V^{*}_{\mathfrak{g}}(\Lambda)
=\Hom_{\mathfrak{g}}(V_{\lambda_n}\otimes\cdots\otimes V_{\lambda_1},\C). \tag{6.29}
\end{equation}
向量空间 $V^{*}_{\mathfrak{g}}(\Lambda)$ 分解为
\begin{equation}
V^{*}_{\mathfrak{g}}(\Lambda)=
\sum_{\vec{\mu}\in(P_\ell)^{n-1}}V^{*}_{\mathfrak{g}}(\Lambda)_{\vec{\mu}},
\qquad
\vec{\mu}=(\mu_{n-1},\mu_{n-2},\dots,\mu_1)\in(P_\ell)^{n-1},
\tag{6.30}
\end{equation}
其间有自然同构
\[
\begin{aligned}
C_{\vec{\mu}}:V^{*}_{\mathfrak{g}}(\Lambda)_{\vec{\mu}}
&\simeq\Hom_{\mathfrak{g}}(V_{\lambda_n}\otimes V_{\mu_{n-1}},V_0)
 \otimes\Hom_{\mathfrak{g}}(V_{\lambda_{n-1}}\otimes V_{\mu_{n-2}},
 V_{\mu_{n-1}})\otimes\\
&\qquad\cdots\otimes\Hom_{\mathfrak{g}}(V_{\lambda_2}\otimes V_{\mu_1},
 V_{\mu_2})\otimes\Hom_{\mathfrak{g}}(V_{\lambda_1}\otimes V_0,V_{\mu_1}).
\end{aligned}
\]
这里线性映射 $C_{\vec{\mu}}$ 定义为
\[
\begin{aligned}
&C_{\vec{\mu}}(\varphi_n\otimes\varphi_{n-1}\otimes\cdots\otimes\varphi_1)
 (u_n\otimes\cdots\otimes u_1)\\
&\quad=\langle0|\varphi_n\bigl(u_n\otimes\varphi_{n-1}\bigl(u_{n-1}\otimes
 \cdots\varphi_2(u_2\otimes\varphi_1(u_1\otimes|0\rangle))\cdots\bigr)\bigr),
\end{aligned}
\]
其中
$\varphi_i\in\Hom_{\mathfrak{g}}(V_{\lambda_i}\otimes V_{\mu_{i-1}},V_{\mu_i})
=\Hom_{\mathfrak{g}}(V_{\lambda_i},\Hom(V_{\mu_{i-1}},V_{\mu_i}))$，
$1\le i\le n$，$\mu_n=\mu_0=0$，且
$u_n\otimes\cdots\otimes u_1\in V_\Lambda
=V_{\lambda_n}\otimes\cdots\otimes V_{\lambda_1}$。
集合 $\vec{\mu}$ 给出一列顶点
\[
\mathbf{v}_n(\vec{\mu})=\binom{\lambda_n}{0\ \mu_{n-1}},\ \dots,\
\mathbf{v}_i(\vec{\mu})=\binom{\lambda_i}{\mu_i\ \mu_{i-1}},\ \dots,\
\mathbf{v}_1(\vec{\mu})=\binom{\lambda_1}{\mu_1\ 0},
\]
并由此定义 $V^{*}_{\mathfrak{g}}(\Lambda)$ 的子空间 $\cV(\Lambda)$：
\[
\cV(\Lambda)=\sum_{\vec{\mu}\in(P_\ell)^{n-1}}\cV_{\vec{\mu}},
\]
其中
\[
\cV_{\vec{\mu}}=C^{-1}_{\vec{\mu}}\Bigl(\cV(\mathbf{v}_n(\vec{\mu}))
\otimes\cdots\otimes\cV(\mathbf{v}_i(\vec{\mu}))\otimes\cdots
\otimes\cV(\mathbf{v}_1(\vec{\mu}))\Bigr).
\]

% [Tutor Note: $\ell C$-约束就是 (6.28)：$L_i=\ell-(\lambda_i,\theta)+1$ 是
%  「第 $i$ 个顶点处可积截断的长度」，而 $\bar m_i$ 在三权之外的所有方向取分量、
%  总和恰为 $L_i$。$C_{\vec{\mu}}$ 的公式把 $\Hom_{\mathfrak{g}}$ 的元素按
%  中间权 $\vec{\mu}=(\mu_{n-1},\dots,\mu_1)$ 逐层拆开：第 $i$ 层是
%  $\Hom_{\mathfrak{g}}(V_{\lambda_i}\otimes V_{\mu_{i-1}},V_{\mu_i})$，
%  从 $V_{\mu_0}=V_0=\C$ 出发逐层收缩，最后回到 $V_{\mu_n}=V_0$。
%  原文 p.146]

\figomitn{与 p.144 末相同的顶点算子链符号图：水平线上立着
$\lambda_n,\lambda_{n-1},\dots,\lambda_1$ 诸顶点，相邻顶点之间标
$\mu_{n-1},\mu_{n-2},\dots,\mu_1$，两端标 $0$。}

% [原文 p.147]

现在，对每个非零的
$\varphi=\varphi_n\otimes\cdots\otimes\varphi_1\in\cV_{\vec{\mu}}$，
可以对应一个 $n$-点函数
\[
\Phi_\varphi(z_n,\dots,z_1)
=\langle\Phi_{\varphi_n}(z_n)\cdots\Phi_{\varphi_1}(z_1)\rangle .
\]
于是由定理 6.17 有下述定理。

\begin{theorem}[6.18]
\begin{enumerate}[label=\arabic*.]
\item 对每个非零的 $\varphi\in\cV^{*}_{\mathfrak{g}}(\Lambda)_{\vec{\mu}}$，
形式 Laurent 级数 $\Phi_\varphi(z_n,\dots,z_1)$ 在区域
\[
\mathcal{R}_n=\{(z_n,\dots,z_1)\in\C^n\mid
|z_n|>|z_{n-1}|>\cdots>|z_1|\}
\]
内绝对收敛，并解析延拓为 $X_n$ 上的多值全纯函数，其中
\[
X_n=\C^n\setminus\bigcup_{i<j}\Delta_{ij},\qquad
\Delta_{ij}=\{(z_n,\dots,z_1)\in\C^n\mid z_i=z_j\}.
\]
\item 带 $\ell C$-约束 (6.28) 的 KZ 方程 (6.27) 的解空间同构于
$\cV(\Lambda)$。
\end{enumerate}
\end{theorem}

% [Tutor Note: 定理 6.18 的两条合起来说：$n$-点函数全体在
%  $\mathcal{R}_n=\{|z_n|>\cdots>|z_1|\}$ 内绝对收敛（形式级数合法），
%  且它们张成带 $\ell C$-约束的 KZ 方程的全部解，故解空间同构于
%  $\cV(\Lambda)$。沿 $X_n$ 中的闭路解析延拓即得单值性表示
%  $\pi_1(X_n)\to GL(\cV(\Lambda))$，这是 Drinfeld–Kohno 定理的出发点。
%  原文 p.147]

由于 KZ 方程的解是多值函数，计算它们的单值性（monodromy），即空间 $X_n$
的基本群到带 $\ell C$-约束的 KZ 方程解空间的表示，是一个有趣的问题。

先考虑四点情形。固定四元组
\[
\Lambda=(\lambda_4,\lambda_3,\lambda_2,\lambda_1).
\]
对每个 $\mu\in P_\ell$ 置
\[
\mathbf{v}_2(\mu)=\binom{\lambda_3}{\lambda_4^{\dagger}\ \mu},\qquad
\mathbf{v}_1(\mu)=\binom{\lambda_2}{\mu\ \lambda_1}.
\]
再置
\[
\Delta_4(\Lambda)=\widehat{\Delta}(\mathbf{v}_2)+\widehat{\Delta}(\mathbf{v}_1)
=\Delta_{\lambda_1}+\Delta_{\lambda_2}+\Delta_{\lambda_3}-\Delta_{\lambda_4}.
\]

% [Tutor Note: 上式右边可由 (6.22) 直接验证：按定义 $\mathbf{v}_2$ 的输出端是
%  $\lambda_4^{\dagger}$、输入端是 $\mu$，故
%  $\widehat{\Delta}(\mathbf{v}_2)=\Delta_{\lambda_3}+\Delta_\mu
%  -\Delta_{\lambda_4^{\dagger}}$；$\mathbf{v}_1$ 的输出端是 $\mu$、输入端是
%  $\lambda_1$，故 $\widehat{\Delta}(\mathbf{v}_1)=\Delta_{\lambda_2}
%  +\Delta_{\lambda_1}-\Delta_\mu$。两式相加 $\Delta_\mu$ 相消，再用对偶模的
%  共形维数不变 $\Delta_{\lambda^{\dagger}}=\Delta_\lambda$，即得印出的等式。
%  原文 p.147]

\figomitn{四顶点链的符号图：水平线上方立着 $\lambda_3$、$\lambda_2$，
下方自左至右标 $\lambda_4^{\dagger}$、$\mu$、$\lambda_1$。}

于是 KZ 方程 $KZ_2(\Lambda)$ 与 $\ell C$-约束 $\ell C_2(\Lambda)$
如下给出。

% [原文 p.148]

\[
\begin{aligned}
KZ_2(\Lambda)\qquad
&\left(\kappa\frac{\partial}{\partial w}-\frac{\Omega_{13}}{w}
-\frac{\Omega_{23}}{w-z}\right)\Psi(w,z)\\
&=\left(\kappa\frac{\partial}{\partial z}-\frac{\Omega_{12}}{w}
-\frac{\Omega_{23}}{z-w}\right)\Psi(w,z)=0,
\end{aligned}
\]
\[
\begin{aligned}
\ell C_2(\Lambda)\qquad
&\sum_{m=0}^{L_1}\binom{L_1}{m}w^{-m}z^{m-L_1}
 \Psi(w,z)(u_4,X_\theta^m u_3,X_\theta^{L_1-m}u_2,|\lambda_1\rangle)=0,\\
&\sum_{m=0}^{L_2}\binom{L_2}{m}(w-z)^{-m}(-z)^{m-L_2}
 \Psi(w,z)(u_4,X_\theta^m u_3,|\lambda_2\rangle,
 X_\theta^{L_2-m}u_1)=0,\\
&\sum_{m=0}^{L_3}\binom{L_2}{m}(w-z)^{-m}(-w)^{m-L_3}
 \Psi(w,z)(u_4,|\lambda_3\rangle,X_\theta^m u_2,
 X_\theta^{L_3-m}u_1)=0,\\
&\sum_{|\bar{m}|=L_4}\binom{L_4}{\bar{m}}w^{-m}z^{m-L_1}
 \Psi(w,z)(|\lambda_4^{\dagger}\rangle,X_\theta^{m_3}u_3,
 X_\theta^{m_2}u_2,X_\theta^{m_1}u_1)=0,
\end{aligned}
\]
其中 $L_i=\ell-(\lambda_i,\theta)+1$，$i=1,\dots,4$，且
$\bar{m}=(m_3,m_2,m_1)\in(\Z_{\ge0})^3$。

% [Tutor Note: 原书 $\ell C_2(\Lambda)$ 第三行把 $\binom{L_3}{m}$ 印成
%  $\binom{L_2}{m}$，第四行把 $w^{-\bar m}$ 印成 $w^{-m}$，且上标
%  $z^{m-L_1}$ 中的 $L_1$ 未随行改动；$KZ_2(\Lambda)$ 第一行第二项亦印作
%  $\Omega_{12}/w$（按 (6.27) 的形状应为 $\Omega_{12}/(w-z)$ 型的差商）。
%  这些都是原文排印疏漏，照原文排，未改。原文 p.148]
%  同类照排原文的印刷问题（p.149）：(i) $\ell C_1(\Lambda)$ 第三行求和上限为 $L_3$，
%  而二项式系数原书印 $\binom{L_2}{m}$，疑应为 $\binom{L_3}{m}$；(ii) $(T\varphi)$ 的
%  定义式左右两端各多印一个 $u_4$（$\varphi$ 只吃 4 个变量）。均未改，见 ERRATA。

于是对任意 $\varphi_2\in\cV(\mathbf{v}_2(\vec{\mu}))$ 与
$\varphi_1\in\cV(\mathbf{v}_1(\vec{\mu}))$ 置
\[
\Psi_{\varphi_2\otimes\varphi_1}(w,z)
=\langle\Phi_{\varphi_2}\Phi_{\varphi_1}\rangle .
\]
则由定理 6.18，带 $\ell C_2(\Lambda)$ 的 $KZ_2(\Lambda)$ 的解空间由所有
$\Psi_{\varphi_2\otimes\varphi_1}(w,z)$
（$\varphi_2\in\cV(\mathbf{v}_2(\vec{\mu}))$，
$\varphi_1\in\cV(\mathbf{v}_1(\vec{\mu}))$）张成。

我们引入 $X_2$ 的子集
\[
I_2=\{(w,z)\in\R^2\mid w>z>0\},\qquad
J_2=\{(z,w)\in\R^2\mid z>w>0\},
\]
并考虑 $X_2$ 中的道路 $\gamma(t)=(\eta(t),\zeta(t))$，$t\in[0,1]$，其中
\[
\eta(t)=\frac{w+z}{2}+\ee^{\pi\sqrt{-1}\,t}\frac{w-z}{2},\qquad
\zeta(t)=\frac{w+z}{2}-\ee^{\pi\sqrt{-1}\,t}\frac{w-z}{2}
\]
（$(w,z)\in I_2$）。于是 $X_2$ 中的道路 $\gamma(t)$ 是从点
$(w,z)\in I_2$ 到点 $(z,w)\in J_2$ 的道路。用
$\Psi^{\gamma}_{\varphi_1\otimes\varphi_2}(z,w)$ 表示沿道路 $\gamma(t)$
解析延拓所得的函数。

\figomit{6.1}{道路 $\gamma(t)$。}

对 $\Lambda=(\lambda_4,\lambda_3,\lambda_2,\lambda_1)$ 置
\[
T\Lambda=(\lambda_4,\lambda_2,\lambda_3,\lambda_1).
\]

% [原文 p.149]

并引入 $\mathfrak{g}$-同构
\[
T:V^{*}(\Lambda)=\Hom(V_{\lambda_4}\otimes V_{\lambda_3}\otimes
V_{\lambda_2}\otimes V_{\lambda_1},\C)\to V^{*}(T\Lambda),
\]
它由
\[
(T\varphi)(u_4\otimes u_4\otimes u_2\otimes u_3\otimes u_1)
=\varphi(u_4\otimes u_4\otimes u_3\otimes u_2\otimes u_1)
\]
（$\varphi\in V^{*}(\Lambda)$）定义。则
\[
T(V^{*}_{\mathfrak{g}}(\Lambda))=V^{*}_{\mathfrak{g}}(T\Lambda),\qquad
T(\cV(\Lambda))=\cV(\Lambda).
\]

于是 $\Psi^{\gamma}_{\varphi_1\otimes\varphi_2}(z,w)$ 满足带 $\ell C$-约束
$\ell C_2(T\Lambda)$ 的 KZ 方程 $KZ_2(T\Lambda)$。由定理 6.18，存在线性同构
\[
C(\Lambda)=C_\gamma(\Lambda):\cV(\Lambda)\simeq\cV(T\Lambda).
\]

下面更仔细地研究同构 $C(\Lambda)$。为此引入新变量 $\zeta=z/w$。
则 $V^{*}_{\mathfrak{g}}(\Lambda)$ 值函数
\[
z^{\Delta_4(\Lambda)}\Psi_{\varphi_2\otimes\varphi_1}(w,\zeta w)
\]
只依赖于变量 $\zeta$。因此置
\[
\Psi_{\varphi_2\otimes\varphi_1}(\zeta)
=z^{\Delta_4(\Lambda)}\Psi_{\varphi_2\otimes\varphi_1}(w,\zeta w).
\]
则 $\Psi_{\varphi_2\otimes\varphi_1}(\zeta)$ 称为约化四点函数
（reduced $4$-point function），并满足联立方程组 $KZ_1(\Lambda)$ 与
$\ell C_1(\Lambda)$。

\[
\begin{aligned}
KZ_1(\Lambda)\qquad
&\left(\kappa\frac{\dd}{\dd\zeta}
-\frac{\Omega_{12}+\kappa\Delta_4(\Lambda)}{w}
-\frac{\Omega_{23}}{\zeta-1}\right)\Psi(\zeta)=0,
\end{aligned}
\]
\[
\begin{aligned}
\ell C_1(\Lambda)\qquad
&\sum_{m=0}^{L_1}\binom{L_1}{m}\zeta^m
 \Psi(\zeta)(u_4,X_\theta^m u_3,X_\theta^{L_1-m}u_2,|\lambda_1\rangle)=0,\\
&\sum_{m=0}^{L_2}\binom{L_2}{m}\left(\frac{\zeta}{\zeta-1}\right)^m
 \Psi(\zeta)(u_4,X_\theta^m u_3,|\lambda_2\rangle,
 X_\theta^{L_2-m}u_1)=0,\\
&\sum_{m=0}^{L_3}\binom{L_2}{m}\left(\frac{1}{\zeta-1}\right)^m
 \Psi(\zeta)(u_4,|\lambda_3\rangle,X_\theta^m u_2,
 X_\theta^{L_3-m}u_1)=0,\\
&\sum_{|\bar{m}|=L_4}\binom{L_4}{\bar{m}}w^{-m}z^{m-L_1}
 \Psi(\zeta)(|\lambda_4\rangle,X_\theta^{m_3}u_3,
 X_\theta^{m_2}u_2,X_\theta^{m_1}u_1)=0,
\end{aligned}
\]
其中 $L_i=\ell-(\lambda_i,\theta)+1$，$i=1,\dots,4$，且
$\bar{m}=(m_3,m_2,m_1)\in(\Z_{\ge0})^3$。

\begin{proposition}[6.19]
联立方程组 $KZ_1(\Lambda)$ 与 $\ell C_1(\Lambda)$ 的解空间是
\[
\{\Psi_{\varphi_2\otimes\varphi_1}(\zeta)\mid
\varphi_i\in\cV(\mathbf{v}_i(\mu)),\ i=1,2,\ \mu\in P_\ell\}.
\]
\end{proposition}

\noindent\textit{证明.}\quad 由于
\[
\Psi_{\varphi_2\otimes\varphi_1}(w,z)
=z^{-\Delta_4(\Lambda)}\Psi_{\varphi_2\otimes\varphi_1}
\left(\frac{w}{z}\right)
\]
是联立方程组 $KZ_2(\Lambda)$ 与 $\ell C_2(\Lambda)$ 的解，命题是定理 6.18
的推论。\hfill$\square$

% [原文 p.150]

方程组 $KZ_1(\Lambda)$ 是 Fuchs 型微分方程，正则奇异点只在 $0$、$1$ 与
$\infty$ 处。在 KZ 方程 $KZ_1(\Lambda)$ 的一个奇异点的邻域内，我们将构造
联立方程组 $KZ_1(\Lambda)$ 与 $\ell C_1(\Lambda)$ 的解空间的一组规范基。置
\begin{equation}
V^{*}_{\mathfrak{g}}(\Lambda)^{12}_{\mu}
=\Hom_{\mathfrak{g}}(V_\mu\otimes V_{\lambda_3},V_{\lambda_4^{\dagger}})
\otimes\Hom_{\mathfrak{g}}(V_{\lambda_2}\otimes V_{\lambda_1},V_\mu),
\tag{6.31}
\end{equation}
\begin{equation}
V^{*}_{\mathfrak{g}}(\Lambda)^{13}_{\mu}
=\Hom_{\mathfrak{g}}(V_\mu\otimes V_{\lambda_2},V_{\lambda_4^{\dagger}})
\otimes\Hom_{\mathfrak{g}}(V_{\lambda_3}\otimes V_{\lambda_1},V_\mu),
\tag{6.32}
\end{equation}
\begin{equation}
V^{*}_{\mathfrak{g}}(\Lambda)^{23}_{\mu}
=\Hom_{\mathfrak{g}}(V_\mu\otimes V_{\lambda_1},V_{\lambda_4^{\dagger}})
\otimes\Hom_{\mathfrak{g}}(V_{\lambda_3}\otimes V_{\lambda_2},V_\mu).
\tag{6.33}
\end{equation}
于是有典范同构
\[
C^{12}_{\Lambda}:\bigoplus_{\mu\in P_\ell}
V^{*}_{\mathfrak{g}}(\Lambda)^{12}_{\mu}\simeq
V^{*}_{\mathfrak{g}}(\Lambda),
\]
\[
C^{13}_{\Lambda}:\bigoplus_{\mu\in P_\ell}
V^{*}_{\mathfrak{g}}(\Lambda)^{13}_{\mu}\simeq
V^{*}_{\mathfrak{g}}(\Lambda),\qquad
C^{23}_{\Lambda}:\bigoplus_{\mu\in P_\ell}
V^{*}_{\mathfrak{g}}(\Lambda)^{23}_{\mu}\simeq
V^{*}_{\mathfrak{g}}(\Lambda).
\]

% [Tutor Note: 上标 $12$、$13$、$23$ 表示四点函数在哪个通道里收缩：$(12)(34)$、
%  $(13)(24)$、$(23)(14)$ 三种配对，对应 $z_1$ 与 $z_2$ 的「结合顺序」。
%  相应的单值性（由 $C_{\Lambda}^{12}$ 等给出）就是量子群 $U_q(\mathfrak{g})$
%  的 $R$-矩阵作用；三个通道分别对应奇异点 $0$、$1$、$\infty$ 的局部解。
%  原文 p.150]

为简单起见，我们还使用记号：
\[
V^{*}_{\mathfrak{g}}(\Lambda)^{(0)}_{\mu}
=V^{*}_{\mathfrak{g}}(\Lambda)^{12}_{\mu},\quad
V^{*}_{\mathfrak{g}}(\Lambda)^{(1)}_{\mu}
=V^{*}_{\mathfrak{g}}(\Lambda)^{23}_{\mu},\quad
V^{*}_{\mathfrak{g}}(\Lambda)^{(\infty)}_{\mu}
=V^{*}_{\mathfrak{g}}(\Lambda)^{13}_{\mu}.
\]
现在取 $V^{*}_{\mathfrak{g}}(\Lambda)^{(i)}_{\mu}$ 的基
$\{U^{(i)}_{\mu,k}(\Lambda)\}$，$i=0,1,\infty$。它们构成
$V^{*}_{\mathfrak{g}}(\Lambda)$ 的基，并使算子 $\Omega_{12}$、$\Omega_{23}$
与 $\Omega_{13}$ 对角化。确切地说有
\[
\begin{aligned}
\Omega_{12}U^{(0)}_{\mu,k}
&=\kappa(\gamma^{(0)}_{\mu}-\Delta_4(\Lambda))U^{(0)}_{\mu,k},\\
\Omega_{23}U^{(1)}_{\mu,k}
&=\kappa\gamma^{(1)}_{\mu}U^{(1)}_{\mu,k},\\
\Omega_{23}U^{(\infty)}_{\mu,k}
&=\kappa\gamma^{(\infty)}_{\mu}U^{(\infty)}_{\mu,k},
\end{aligned}
\]
其中
\[
\begin{aligned}
\gamma^{(0)}_{\mu}&=\Delta_\mu-\Delta_{\lambda_3}-\Delta_{\lambda_4}
=\widehat{\Delta}(\mathbf{v}_2(\mu)),\\
\gamma^{(1)}_{\mu}&=\Delta_\mu-\Delta_{\lambda_2}-\Delta_{\lambda_3},\\
\gamma^{(\infty)}_{\mu}&=\Delta_\mu-\Delta_{\lambda_1}-\Delta_{\lambda_3}.
\end{aligned}
\]

在奇异点 $1$ 与 $\infty$ 处取坐标为
\[
\xi=\zeta-1,\qquad \eta=1/\zeta .
\]
则关于这些坐标 $\xi$ 与 $\eta$，KZ 方程 $KZ_1(\Lambda)$ 改写为
\[
\begin{aligned}
KZ_1(\Lambda)_1\qquad
&\left(\frac{\dd}{\dd\xi}-\frac{\Omega_{23}}{\xi}
-\frac{\Omega_{12}+\kappa\Delta_4(\Lambda)}{\xi-1}\right)\Psi(1-\xi)=0,\\
KZ_1(\Lambda)_\infty\qquad
&\left(\frac{\dd}{\dd\eta}-\frac{\Omega_{13}}{\eta}
-\frac{\Omega_{23}}{\eta-1}\right)\Psi\left(\frac{1}{\eta}\right)=0,
\end{aligned}
\]
这里 $\Psi(\zeta)$ 是 $X_1=\C^*$ 上的 $V^{*}_{\mathfrak{g}}(\Lambda)$ 值函数。

% [原文 p.151]

于是由命题 6.12，我们可以选取 KZ 方程 $KZ_1(\Lambda)$ 解空间的三组不同基
\[
\left\{\Psi^{(i)}_{U^{(i)}_{\mu,k}(\Lambda)}(\zeta^{(i)})\right\}
\qquad(\zeta^{(0)}=\zeta,\ \zeta^{(1)}=\xi,\ \zeta^{\infty}=\eta),
\]
它们分别在各奇异点 $i=0,1,\infty$ 处归一化如下：
\[
\begin{aligned}
\Psi^{(0)}_{U^{(0)}_{\mu,k}(\Lambda)}(\zeta)
&=\zeta^{\gamma^{(0)}_{\mu}}
 \left\{C^{12}_{\Lambda}(U^{(0)}_{\mu,k}(\Lambda))+O(\zeta)\right\},\\
\Psi^{(1)}_{U^{(1)}_{\mu,k}(\Lambda)}(\xi)
&=\xi^{\gamma^{(1)}_{\mu}}
 \left\{C^{23}_{\Lambda}(U^{(1)}_{\mu,k}(\Lambda))+O(\xi)\right\},\\
\Psi^{(\infty)}_{U^{(\infty)}_{\mu,k}(\Lambda)}(\eta)
&=\eta^{\gamma^{(\infty)}_{\mu}}
 \left\{C^{13}_{\Lambda}(U^{(\infty)}_{\mu,k}(\Lambda))+O(\eta)\right\}.
\end{aligned}
\]

% [Tutor Note: 这三行是「在奇点处的归一化基」：在 $\zeta=0$ 附近解的行为是
%  $\zeta^{\gamma^{(0)}_{\mu}}$ 乘一个解析因子，而首项系数恰好是
%  $C^{12}_{\Lambda}$ 给出的配对元素。三个奇点给出的三组基之间的过渡矩阵
%  就是下面命题 6.20 的「连接同构」，它是量子群 $R$-矩阵的解析实现。
%  原文 p.151]

现在可以证明下述命题。

\begin{proposition}[6.20]
对任意四元组 $\Lambda\in(P_\ell)^4$，同构 $C(\Lambda)$ 由联立方程组
$KZ_1(\Lambda)$ 与 $\ell C_1(\Lambda)$ 的解空间在 $0$ 处正则化到在
$\infty$ 处正则化的连接同构给出。
\end{proposition}

\noindent\textit{证明.}\quad 对每个 $U\in\cV(\Lambda)^{12}_{\mu}$，
$KZ_2(\Lambda)$ 与 $\ell C_2(\Lambda)$ 的解 $\Psi_U(w,z)$ 在区域
$\mathcal{R}_2$ 内绝对收敛。用 $\Psi^{\gamma}_U(z,w)$ 表示它沿道路
$\gamma(t)$ 的解析延拓。则 $T\Psi^{\gamma}_U(z,w)$ 是 $KZ_2(T\Lambda)$ 与
$\ell C_2(T\Lambda)$ 的解，且等于 $\Psi_{\bar U}(w,z)$，其中
$\bar U=C(\Lambda)(U)\in\sum_{\mu\in P_\ell}\cV(T\Lambda)^{12}_{\mu}$。
于是有
\[
T\Psi_U\left(\frac{1}{\zeta}\right)
=\zeta^{-\Delta_4(\Lambda)}\Psi_{\bar U}(\zeta),
\]
其中 $\Psi_U(1/\zeta)$ 是 $\Psi_U(\zeta)$ 沿从 $0$ 到 $\infty$ 的道路
$\bar{\gamma}(t)$ 解析延拓所得者。由于有
\[
\gamma^{(\infty)}_{\mu}(\Lambda)+\Delta_4(\Lambda)
=\gamma^{(0)}_{\mu}(T\Lambda),
\]
且连接同构保持子空间 $\cV(\Lambda)$，我们便得到所要的结果。
\hfill$\square$

下一节我们将对 $\mathfrak{g}=\mathfrak{sl}(n+1,\C)$ 显式地计算连接同构。
